\chapter{Correction Ledger}
\label{app:corrections}

This appendix preserves a visible record of corrections made during the final
evidence audit. It is intentionally included rather than silently replacing
earlier claims: the first evaluation was a single-run demonstration, whereas
the final campaign consists of repeated replay runs and inferential analysis.

\small
\begin{longtable}{p{2.7cm}p{3.4cm}p{4.7cm}}
\caption{Correction ledger: initial claim, corrected statement, and evidence.}\label{tab:correction-ledger}\\
\toprule
\textbf{Location} & \textbf{Earlier claim} & \textbf{Corrected statement and evidence} \\
\midrule
\endfirsthead
\toprule
\textbf{Location} & \textbf{Earlier claim} & \textbf{Corrected statement and evidence} \\
\midrule
\endhead
Abstracts; Chapter 7 & MAS is approximately 36--38\% faster to alert. & The 80.4\% reduction occurs only on \texttt{demo\_site\_01} ($n=5$, $p=0.0625$). Across the seven scenarios having all four modes the reduction is 10.7\% ($p=0.59$), not significant. Source: \texttt{thesis\_stats.json}. \\
Abstracts; Chapter 7 & Auction coordination halves false alerts and restores precision. & Auction has the lowest marginal mean F1 (0.452). After Holm--Bonferroni correction, none of six pairwise F1 comparisons is significant. Sources: \texttt{summary\_agg.csv}; \texttt{thesis\_stats.json}. \\
Chapter 6, earlier interpretation & ``Who verifies'' validates the utility-bid design. & Retracted: no campaign manifest populated \texttt{fov\_overlap}; the overlap factor was the 0.5 default. The bid is therefore not a measured view-utility allocation. Sources: \texttt{camera\_agent.py}; coordinator task construction. \\
Chapter 6, earlier fusion text & Confidence 0.78 demonstrates an audio--video corroboration bonus. & Retracted: the 0.78 alert contains one video event from \texttt{cam\_01}. Four of 240 alerts are truly cross-modal, all at 1.0 because noisy-OR saturates. Source: \texttt{cross\_modal\_fusion\_audit.md}. \\
Chapter 6; Chapter 7 & No hallucinated citations were found in replay. & The campaign never invoked \texttt{--llm}; all 373 explanations used the template. One post-campaign model output was rejected by the guardrail. The claim is a design/property test, not campaign evidence. \\
Chapter 5 & The CLIP anomaly component is presented as working. & Its calibration AUC is 0.308, below chance, consistent with prompt/domain mismatch. Source: \texttt{clip\_anomaly\_calibration\_notes.md}. \\
Chapter 6 setup & Variants differ only in architecture. & They also differ in replay pacing and stream sequencing; fusion and cooldown are wall-clock mechanisms. Source: \texttt{replay.py}. \\
Chapters 1, 5, 7 & Approximately 2,500 lines, seven agents, six tests, and a \texttt{configs/} directory. & 4,376 tracked Python lines (5,188 including working-tree files); six implemented agents plus an operator console; 26 offline tests; zone data and thresholds are scenario manifests, not a \texttt{configs/} directory. \\
Chapters 1, 3, 4 & No Re-ID. & Retained as the evaluated-system claim. A post-campaign appearance-descriptor prototype was withheld; it did not form part of the reported configuration. \\
THESIS\_REPATCH variance row & Identical auction precision values establish stability. & Rejected: Levene tests give $p=0.224$ for \texttt{demo\_site\_01} precision and $p=0.515$ pooled. Five identical observations do not establish lower variance. \\
FA/h figures & False alarms per hour is an operational rate. & Over 9--70 second replay runs with 0--3 false positives, FA/h is a scale artefact (maximum 393.9/h), not a deployment-rate estimate. \\
Chapter 7 & The audit stream is immutable. & It is append-only by convention, not cryptographically immutable. \\
\bottomrule
\end{longtable}
\normalsize

\chapter{Campaign Scenario Manifests}
\label{app:scenarios}

The following are the nine JSON scenario manifests used by the v2 campaign.
They are reproduced directly from the repository so that sensor topology, zone
geometry, and ground-truth intervals can be inspected without transcription.
The three later exploratory manifests (\texttt{loitering\_multizone\_01},
\texttt{wrong\_direction\_01}, and \texttt{zone\_occupancy\_01}) are not part
of the 373-run campaign and are therefore not reported as campaign evidence.

\section{Abandoned-object scenario}
\VerbatimInput[fontsize=\scriptsize,frame=single]{../scenarios/abandoned_object_01.json}
\section{Clock-alarm scenario}
\VerbatimInput[fontsize=\scriptsize,frame=single]{../scenarios/audio_alarm_clock_01.json}
\section{Siren-alarm scenario}
\VerbatimInput[fontsize=\scriptsize,frame=single]{../scenarios/audio_alarm_siren_01.json}
\section{Glass-break scenario}
\VerbatimInput[fontsize=\scriptsize,frame=single]{../scenarios/audio_glass_break_01.json}
\section{Combined audio--video scenario}
\VerbatimInput[fontsize=\scriptsize,frame=single]{../scenarios/combined_audio_video_01.json}
\section{Demonstrator site scenario}
\VerbatimInput[fontsize=\scriptsize,frame=single]{../scenarios/demo_site_01.json}
\section{Fight scenario}
\VerbatimInput[fontsize=\scriptsize,frame=single]{../scenarios/fight_01.json}
\section{Intrusion scenario}
\VerbatimInput[fontsize=\scriptsize,frame=single]{../scenarios/intrusion_01.json}
\section{Loitering true-negative probe}
\VerbatimInput[fontsize=\scriptsize,frame=single]{../scenarios/loitering_01.json}

\chapter{Data Provenance and Licence Notes}
\label{app:datasets}

Table~\ref{tab:data-provenance} transcribes the provenance information in
\texttt{data/clips\_real/manifest.json}. These clips extend the original
demonstrator assets; they are used only for non-commercial academic evaluation.

\small
\begin{longtable}{p{2.8cm}p{3.5cm}p{2.6cm}p{3.3cm}}
\caption{Provenance and licence notes for real evaluation assets.}\label{tab:data-provenance}\\
\toprule
\textbf{Asset} & \textbf{Source} & \textbf{Licence / use note} & \textbf{Role in corpus} \\
\midrule
\endfirsthead
\toprule
\textbf{Asset} & \textbf{Source} & \textbf{Licence / use note} & \textbf{Role in corpus} \\
\midrule
\endhead
\texttt{violence/violent\_1.mp4} & Bianculli et al. \cite{bianculli2020violence} & Free for research and educational purposes & Simulated brawl, about 4.8 s, 1920$\times$1080 at 30 FPS. \\
\texttt{violence/violent\_10.mp4} & Same dataset & Free for research and educational purposes & Available corpus asset; not a separately reported manifest. \\
\texttt{violence/nonviolent\_1.mp4} & Same dataset & Free for research and educational purposes & Deliberately ambiguous non-violent hard negative. \\
\texttt{abandoned\_object/video1.avi} & ABODA \cite{lin2015aboda} & Public research dataset; no explicit per-file licence in its repository & 73-second indoor hallway sequence. Its onset is an approximate, frame-checked annotation because ABODA supplies no frame-level ground truth. \\
\texttt{audio/glass\_breaking\_esc50.wav} & ESC-50 \cite{piczak2015esc50} & CC BY-NC 3.0 & Real five-second glass-breaking recording. \\
\texttt{audio/clock\_alarm\_esc50.wav} & ESC-50 \cite{piczak2015esc50} & CC BY-NC 3.0 & Alarm proxy event. \\
\texttt{audio/siren\_esc50.wav} & ESC-50 \cite{piczak2015esc50} & CC BY-NC 3.0 & Five-second siren hazard event. \\
\bottomrule
\end{longtable}
\normalsize

\chapter{Configuration and Measurement Parameters}
\label{app:config}

The reported scores are outputs of a system with numerous free parameters.
Table~\ref{tab:parameters} makes the principal defaults explicit, including
their implementation provenance. It is a disclosure table, not a tuning study:
apart from the documented CLIP threshold sweep, these values were not selected
on a held-out validation split.

\small
\begin{longtable}{p{3.7cm}p{2.1cm}p{5.9cm}}
\caption{Principal free parameters and source locations.}\label{tab:parameters}\\
\toprule
\textbf{Parameter} & \textbf{Default / campaign value} & \textbf{Implementation provenance} \\
\midrule
\endfirsthead
\toprule
\textbf{Parameter} & \textbf{Default / campaign value} & \textbf{Implementation provenance} \\
\midrule
\endhead
Fusion window & 6 s & \texttt{fusion\_agent.py}, \texttt{FusionAgent.window\_seconds} \\
Video reliability weight & 0.9 & \texttt{fusion\_agent.py}, noisy-OR weighting \\
Audio reliability weight & 0.7 & \texttt{fusion\_agent.py}, noisy-OR weighting \\
Cross-modality bonus & 0.05 & \texttt{fusion\_agent.py}, confidence adjustment \\
Cross-sensor bonus & 0.05 & \texttt{fusion\_agent.py}, confidence adjustment \\
Verification gray zone & [0.35, 0.75) & \texttt{coordinator\_agent.py}, \texttt{needs\_verification()} \\
Bid window & 1.0 s & \texttt{coordinator\_agent.py}, constructor default \\
Verification timeout & 3.0 s & \texttt{coordinator\_agent.py}, constructor default \\
Origin-sensor bid factor & 0.3 & \texttt{camera\_agent.py}, \texttt{\_view\_score()} \\
Busy-camera capacity factor & 0.2 & \texttt{camera\_agent.py}, \texttt{\_view\_score()} \\
FoV overlap fallback & 0.5 & \texttt{camera\_agent.py}: task default; inactive as an information-bearing term in campaign measurement \\
Critical / warning / info thresholds & 0.45 / 0.55 / 0.70 & \texttt{policy\_agent.py}, \texttt{ALERT\_THRESHOLDS} \\
Verification adjustment & +0.15 / -0.20 & \texttt{policy\_agent.py}, verification outcome handling \\
Alert cooldown & 20 s & \texttt{policy\_agent.py}, constructor default \\
Inference cadence & 5 FPS & \texttt{camera\_agent.py}, \texttt{infer\_fps} default \\
Camera confidence floor & 0.35 & \texttt{camera\_agent.py}, detector configuration \\
Verification confidence / size & 0.25 / 960 px & \texttt{camera\_agent.py}, \texttt{\_verify()} \\
CLIP anomaly threshold & 0.55 & \texttt{camera\_agent.py}, constructor default \\
Loitering threshold & 8 s & \texttt{camera\_agent.py}, \texttt{ZoneRuleEngine} default \\
Abandoned-object threshold & 10 s & \texttt{camera\_agent.py}, \texttt{ZoneRuleEngine} default \\
Static-object IoU & 0.6 & \texttt{camera\_agent.py}, \texttt{ZoneRuleEngine} default \\
Metric time tolerance & $\pm5$ s & \texttt{metrics.py}, event matching logic \\
\bottomrule
\end{longtable}
\normalsize

\chapter{Methodology Changes from v1 to v2}
\label{app:methodology}

This table is reproduced from \texttt{results/methodology\_changes.md}. It
records changes that could move a result relative to the preserved pre-YAMNet
baseline, avoiding the appearance that v1 and v2 scores measure an unchanged
pipeline.

\small
\begin{longtable}{p{0.5cm}p{3.2cm}p{3.5cm}p{4.0cm}}
\caption{Documented changes between the v1 and v2 evaluation campaigns.}\label{tab:methodology-changes}\\
\toprule
\# & \textbf{Change} & \textbf{Why} & \textbf{Expected direction} \\
\midrule
\endfirsthead
\toprule
\# & \textbf{Change} & \textbf{Why} & \textbf{Expected direction} \\
\midrule
\endhead
1 & Local YAMNet integration instead of DSP fallback only & The TensorFlow Hub URL was dead; this blocked typed audio events. & Audio-scenario recall and precision up. \\
2 & Lookback-extended window and max pooling instead of bare one-second mean pooling & Mean pooling diluted brief transients (0.069 versus 0.75 on the same clip). & Audio recall up. \\
3 & One event per type per audio chunk & Multiple matched classes otherwise entered noisy-OR as independent evidence. & Inflated confidence and false positives down. \\
4 & Unified event-family mapping & Metric and fusion family dictionaries had drifted. & Correctly mapped audio events can score as matches. \\
5 & Audio events carry the scenario zone & Without the zone, audio and video could not join the same fusion hypothesis. & Confidence can rise only for time- and zone-aligned evidence. \\
6 & Scenario clock starts after model warm-up & Cold-start loading had been counted in time-to-alert. & Wall seconds and apparent alert delay down. \\
7 & Repetition flags and rep-suffixed outputs & Reruns previously overwrote results. & Enables mean and standard-deviation reporting. \\
8 & Aggregation to \texttt{summary\_agg.csv} & Repeated rows needed a separate aggregation artifact. & No direct scoring effect. \\
9 & Two additional ESC-50 alarm scenarios & Earlier audio scenarios exercised only glass-break ground truth. & Expands class-discrimination coverage. \\
10 & Full nine-scenario, four-mode, five-repetition campaign & Single-run demonstration could not support uncertainty estimates. & Enables paired inferential analysis. \\
\bottomrule
\end{longtable}
\normalsize

\chapter{Message Schemas}
\label{app:schemas}

The implementation uses typed data classes in \texttt{aura\_mas/core/bus.py}.
The abbreviated field inventory below distinguishes transient detections from
durable semantic events and operator-facing alerts.

\begin{table}[H]
\centering
\small
\caption{Core message-schema fields.}\label{tab:schemas}
\begin{tabular}{p{2.3cm}p{9.5cm}}
\toprule
\textbf{Schema} & \textbf{Key fields} \\
\midrule
Detection & \texttt{sensor\_id}, \texttt{frame\_id}, \texttt{timestamp}, object class, confidence, bounding box, and track identifier. \\
Event & \texttt{event\_id}, \texttt{sensor\_id}, \texttt{timestamp}, \texttt{event\_type}, confidence, modality, zone, track identifier, evidence path, and extensible metadata. \\
Alert & \texttt{alert\_id}, timestamp, severity, event type, confidence, zone, sensor list, evidence list, fused-event list, explanation, and status. \\
\bottomrule
\end{tabular}
\end{table}

The topic constants in the same module route detections, events, coordination
tasks, bids, awards, verifications, alerts, and audit records. This explicit
separation makes it possible to exchange structured facts rather than raw
footage across the system boundary.
